
\documentclass[11pt,a4paper]{article}

% --- Margins & Layout ---
\usepackage[margin=1in]{geometry}
\usepackage{microtype}

% --- Math Packages & Fonts ---
\usepackage{amsmath,amssymb,amsthm,amsfonts}
\usepackage{mathtools}
\usepackage{mathrsfs}

% --- Hyperlinks ---
\usepackage[colorlinks=true, linkcolor=blue, citecolor=blue, urlcolor=blue]{hyperref}

% --- Custom Macros ---
\newcommand{\bR}{\mathbb{R}}
\newcommand{\cC}{\mathcal{C}}
\newcommand{\cM}{\mathcal{M}}
\newcommand{\cL}{\mathcal{L}}

% --- Theorem Environments ---
\theoremstyle{plain}
\newtheorem{theorem}{Theorem}[section]
\newtheorem{lemma}[theorem]{Lemma}
\newtheorem{proposition}[theorem]{Proposition}
\newtheorem{corollary}[theorem]{Corollary}

\theoremstyle{definition}
\newtheorem{definition}[theorem]{Definition}
\newtheorem{example}[theorem]{Example}
\newtheorem{remark}[theorem]{Remark}

\title{\textbf{Ana III Zusammenfassung}}
\author{Moritz}

\begin{document}
\maketitle

\section{Measure Theory}
\begin{definition}[outer measure]
    Let $X$ be any set and let $\mu: 2^X \to [0, \infty]$ then $\mu$ is an 
    outer measure if
    \begin{enumerate}
        \item $\mu(\emptyset) = 0$
        \item $\mu(A) \le \sum_k \mu(A_k)$ \quad for every $A \subset \bigcup_k A_k$
    \end{enumerate}
\end{definition}

\begin{definition}[measurability]
    A set $A \subset X$ is considered $\mu$-measurable if for every $B \subset X$ we have 
    the inequality
    \[
        \mu(B) \ge \mu(A \cap B) + \mu (B \backslash A)
    \]
    Now let $\cM(\mu)$ be the set of all $\mu$-measurable sets.
\end{definition}

\begin{theorem}[basic properties of measurable sets]
    Let $A \subset X$ be measurable and let $A_k \subset X$ be a sequence of measurable sets then 
    we have the following properties 
    \begin{enumerate}
        \item $X\setminus A = A^C \in \cM(\mu)$
        \item $\bigcap_k A_k \in \cM(\mu)$ and $\bigcup_k A_k \in \cM(\mu)$
        \item if $A_k$ is a pairwise disjoint sequence then we have $\mu(\bigcup_k A_k) = \sum_k \mu(A_k)$
    \end{enumerate}
\end{theorem}


\begin{definition}[measure space]
    A triple $(X, \cM, \mu)$ is called a measure space if $\mu: \cM \to [0, \infty]$ is a measure 
    and $\cM$ is a $\sigma$-Algebra. 
\end{definition}

\begin{theorem}[monotone convergence for measures]
    Let $(A_n)_n$ be a sequence of measurable sets
    \begin{enumerate}
        \item if $A_1 \subset A_2 \subset \dots$ then we have 
            \[
                \mu(\bigcup_n A_n) = \lim_{n \to \infty} \mu(A_n)
            \]
        \item if $A_1 \supset A_2 \supset \dots$ and $\mu(A_1) < \infty$ then we have 
            \[
                \mu(\bigcap_n A_n) = \lim_{n \to \infty} \mu(A_n)
            \]
    \end{enumerate}
\end{theorem}


\begin{definition}[$\sigma$-finite]
    A measure space $(X, \cM, \mu)$ is $\sigma$-finite if there are sets $E_n \in \cM$ such that 
    \[
        X = \bigcup_n E_n, \quad \mu(E_n) < \infty
    \]
\end{definition}

\subsection{Lebesgue Measure}
\begin{definition}[rectangle]
    Let $a, b \in \bR^n$ such that $a_i \le b_i$ then we define the rectangle with corners $a$ and $b$ by
    \[
        Q(a, b) := (a_1, b_1) \times \cdots \times (a_n, b_n)
    \]
    and we define the volume with 
    \[
        |Q(a, b)| := \prod_{i=1}^n (b_i - a_i)
    \]
\end{definition}


\begin{definition}[lebesgue outer measure]
    Let $E \subset \bR^n$ then we define the $n$ dimensional Lebesgue Outer Measure by 
    \[
        \cL^n(E) := \inf \{ \sum_k |Q_k|: E \subset \bigcup_k Q_k, \text{ where } Q_k \text{ are rectangles} \}
    \]
\end{definition}

\begin{theorem}[properties of the lebesgue measure]
    The Lebesgue measure has the following properties
    \begin{enumerate}
        \item $\cL^n$ is Borel regular and $\sigma$-finite
        \item $\cL^n(Q(0, 1)) = 1$
        \item $\cL^n(A + t) = \cL^n(A)$ where $t \in \bR^n$
        \item $\cL^n(\Lambda(A)) = \det (\Lambda) \cL^n(A)$ 
            with $\Lambda = \operatorname{diag}(\lambda_1, \dots, \lambda_n)$ and $\lambda_k \ge 0$ for all $k$
    \end{enumerate}
\end{theorem}


\end{document}
